\documentclass[12pt,letterpaper, onecolumn]{exam}
\usepackage{amsmath}
\usepackage{amssymb}
\usepackage[lmargin=71pt, tmargin=1.2in]{geometry}
\usepackage{xcolor}
\usepackage[brazil]{babel}
\usepackage{natbib}
\usepackage{enumitem}
\usepackage{booktabs}
\usepackage[hidelinks]{hyperref}
\urlstyle{same}


\renewenvironment{solution}
  {\par\noindent\textbf{R:}\enspace\color{blue}}
  {\par\color{black}}

  
% \chead{\hline} % Un-comment to draw line below header
\thispagestyle{empty}   %For removing header/footer from page 1

\begin{document}


\begingroup  
    \centering
    \LARGE Lista 6 - Econometria I - 2026\\[0.5em]
    %\large \today\\[0.5em]
    \large Prof: André Luiz Pereira Mancha \par
    \large Monitor: Tuffy Licciardi Issa  \par
    %\large Roll Number\par
    %\large Class/Batch/Section\par
\endgroup
\rule{\textwidth}{0.4pt}
\pointsdroppedatright   %Self-explanatory
\printanswers
\renewcommand{\solutiontitle}{\noindent\textbf{Ans:}\enspace}   %Replace "Ans:" with starting keyword in solution box
\newcommand{\qstar}{\hspace{-0.6em}\textsuperscript{*}\hspace{0.2em}}

\begin{questions}

    \question\qstar (4.11 \citep{wooldridge2016}) A tabela seguinte foi criada ao usar os dados do arquivo \texttt{CEOSAL2}, onde erros padrão estão entre parênteses abaixo dos coeficientes:

\begin{table}[h!]
\centering
\begin{tabular}{lccc}
\toprule
\hline
\\[-1.0ex]
& \multicolumn{3}{c}{Variável dependente: $\log(salary)$} \\[1.5ex]
\cline{2-4}\\[-1.0ex]
Variáveis Independentes & (1) & (2) & (3) \\[1.5ex]
\hline
$\log(sales)$   & 0.224 & 0.158 & 0.188 \\
                & (0.027) & (0.040) & (0.040) \\[0.2cm]

$\log(mktval)$  & --- & 0.101 & 0.097 \\
                & --- & (0.050) & (0.049) \\[0.2cm]

$profmarg$      & --- & -0.0023 & -0.0022 \\
                & --- & (0.0022) & (0.0021) \\[0.2cm]

$ceoten$        & --- & --- & 0.0174 \\
                & --- & --- & (0.0057) \\[0.2cm]

$comten$        & --- & --- & -0.0092 \\
                & --- & --- & (0.0033) \\[0.2cm]

intercepto      & 4.94 & 4.62 & 4.57 \\
                & (0.20) & (0.25) & (0.25) \\[0.2cm]

Observações     & 177 & 177 & 177 \\
R-quadrado      & 0.281 & 0.304 & 0.353 \\
\midrule
\bottomrule
\end{tabular}
\end{table}

A variável \texttt{mktval} é o valor de mercado da empresa, \texttt{profmarg} é o lucro como porcentagem das vendas, \texttt{ceoten} corresponde aos anos de trabalho como diretor-executivo na atual empresa, e \texttt{comten} é o total de anos na empresa.

\begin{enumerate}[label=(\roman*)]
    \item Comente os efeitos de \texttt{profmarg} sobre o salário dos diretores-executivos.
    \item O valor de mercado tem um efeito significante? Explique.
    \item Interprete os coeficientes de \texttt{ceoten} e \texttt{comten}. As variáveis são estatisticamente significantes?
    \item A evidência indica que, permanecendo muito tempo na empresa, mantendo fixos os outros fatores, está associada a salários mais baixos?
\end{enumerate}

    \
\
\


\question (8 (ANPEC, 2012), \citep{pereda_denisard}) Usando uma base de dados que tem informações de 65.535 trabalhadores, queremos verificar se existe desigualdade salarial entre os setores da economia. Considere que a economia está dividida em quatro setores: indústria, comércio, serviços e construção. Cada um dos trabalhadores está em um dos quatro setores e eles são mutuamente exclusivos. Seja \(Y_i\) o salário mensal do trabalhador \(i\) e definimos para cada setor uma variável binária que é igual a 1 se o trabalhador pertence ao setor e 0 caso contrário. Estimando um modelo linear de regressão, obtemos o seguinte resultado:

\[
\hat{Y}_i = 4{,}00 + 0{,}12\,educ_i + 0{,}03\,idade_i + 0{,}40\,homem_i - 0{,}05\,DI_i - 0{,}15\,DC_i - 0{,}25\,DCons_i
\] 
\vspace{-0.8cm}
\[
\quad (0{,}02)\quad (0{,}008)\quad (0{,}0001)\qquad (0{,}0005)\qquad (0{,}0011)\qquad (0{,}003)\qquad (0{,}005)
\]
\[
R^2 = 0{,}83
\]

Em que \textit{educ} representa o número de anos de estudo de cada trabalhador, \textit{idade} é medida em anos, \textit{homem} é uma variável binária que assume valor igual a 1 se é homem e 0 caso contrário, \(DI\) representa a dummy para indústria, \(DC\) para comércio e \(DCons\) para construção. Entre parênteses encontra-se o erro-padrão de cada estimativa.

Baseado nas informações anteriores, julgue as seguintes afirmativas:

\begin{enumerate}[label=(\arabic*)]
    \item Com base nos resultados anteriores é possível rejeitar ao nível de 5\% de significância a hipótese nula de que o salário do setor da indústria é igual ao salário do setor de serviços para trabalhadores com o mesmo nível educacional, a mesma idade e do mesmo sexo. A hipótese alternativa é que os salários nestes setores sejam diferentes.
    
    \item Com base nos resultados anteriores é possível rejeitar ao nível de 5\% de significância a hipótese nula de que o salário no setor de construção é igual ao salário no setor de comércio, mantendo educação, idade e sexo fixos. A hipótese alternativa é que os salários nesses setores sejam diferentes.
    
    \item Com base nos resultados anteriores, é possível rejeitar ao nível de 5\% de significância a hipótese nula de que os salários nos 4 setores da economia são iguais, mantendo constante educação, idade e sexo.
    
    \item Os resultados do modelo anterior permitem testar a hipótese nula de que o retorno salarial entre homem e mulher é diferente para cada nível educacional, ao nível de 5\% de significância.
    
    \item Com base nos resultados anteriores, podemos testar a hipótese de que o intercepto do modelo linear de salário em função de educação, idade e setor para homem é diferente do intercepto do mesmo modelo linear de salário para mulher.
\end{enumerate}
\
\


    
\question\qstar (5.1 \citep{wooldridge2016}) No modelo de regressão simples sob os RLM.1 até RLM.4, afirmamos que o estimador de inclinação, $\hat\beta_1$, é consistente com $\beta_1$. Usando $\beta_0 = \overline y - \hat\beta_1\overline x_1$, demonstre que plim $\hat\beta_0 = \beta_0$. [Você precisará usar a consistência de $\hat\beta_1$ e a lei dos grandes números, juntamente com o fato de que $\beta_o = \text{E}[y] - \beta_1\text{E}[x_1]$].


\
\


\question (5.2 \citep{wooldridge2016}) Suponha que o modelo:
\begin{equation}
    pctstck = \beta_0 + \beta_1funds + \beta_2risktol + u
\end{equation} satisfaça as quatro primeiras hipóteses de Gauss-Markov, em que \texttt{pctstck} é a porcentagem da pensão de um trabalhador investida no mercado de ações, \texttt{funds} é o número de fundos mútuos que o trabalhador pode escolher e \texttt{risktol} é alguma medida de tolerância de risco ( \texttt{risktol} maior significa que a pessoa tem uma tolerância maior ao risco). Se  \texttt{funds} e  \texttt{risktol} são positivamente correlacionados, qual é a inconsistência em $\tilde\beta_1$, o coeficiente de inclinação da regressão  \texttt{pctstck} sobre \texttt{funds}?

\
\\


\question\qstar (5.C1 \citep{wooldridge2016})\footnote{Esta é uma questão computacional. Para instalar o \texttt{R} e o \texttt{RStudio} veja o passo a passo aqui: \href{https://rstudio-education.github.io/hopr/starting.html}{https://rstudio-education.github.io/hopr/starting.html}. Caso tenha dúvidas quanto à linguagem, consulte o material disponibilizado no Moodle, e recomendo também o seguinte material de introdução ao \texttt{R}: \href{https://fhnishida.netlify.app/project/rec2301/sec2/}{https://fhnishida.netlify.app/project/rec2301/sec2/} }
Use os dados do arquivo \texttt{WAGE1} para resolver este exercício.

    \begin{enumerate} [label = (\alph*)]
\item Estime a equação
    \[
    wage = \beta_0 + \beta_1 educ + \beta_2 exper + \beta_3 tenure + u
    \]

    Salve os resíduos e trace um histograma.
\item  Repita o item (a), mas com \(\log(wage)\) como a variável dependente.
\item  Você diria que a Hipótese RLM.6 está mais próxima de ser satisfeita pelo modelo nível-nível ou pelo modelo semilog?
    \end{enumerate}
    
\end{questions}
\newpage
\bibliographystyle{apalike}   % or plainnat, econ, etc.
\bibliography{refs}
\end{document}